\subsection{Неупругие удары}
%%%%%%%%%%%%%%%%%%%%%%%%%%%%%%%%%%%%%%%%%%%%%%%%%%%%
%%%%%%%%%%%%%%%%%%%%%%%%%%%%%%%%%%%%%%%%%%%%%%%%%%%%
\z[{\begin{wrapfigure}[6]{r}{0.2\textwidth}%
		\vspace{-10pt}%
		\begin{flushright}%
 			\includegraphics[width=0.2\textwidth]{7.8.1}%
		\end{flushright}%
	\end{wrapfigure}%
	}]%
	{Пуля массы $m$, имеющая начальную скорость $v$, пробивает подвешенный на нити груз той же массы $m$ и застревает во втором таком же. Найдите выделившееся в первом грузе количество теплоты, если во втором грузе выделилось количество теплоты $Q_2$. Временем взаимодействия пули с грузом пренебречь.}
%
{$Q_1 = 2\sqrt{Q_2 m}\left(v - 2\sqrt{\frac{Q_2}{m}}\right)$
}
%%%%%%%%%%%%%%%%%%%%%%%%%%%%%%%%%%%%%%%%%%%%%%%%%%%%	
\z[{\begin{wrapfigure}[5]{r}{0.2\textwidth}%
		\vspace{-10pt}%
		\begin{flushright}%
 			\includegraphics[width=0.2\textwidth]{7.8.2}%
		\end{flushright}%
	\end{wrapfigure}%
	}]%
	{На столе лежит брусок. На лёгкой нити длиной $L$ висит шарик, касаясь бруска. Нить вертикальна (см. рисунок). Масса бруска в 7 раз больше массы шарика. Шарик отклоняют в сторону так, что нить занимает горизонтальное положение, и отпускают. После неупругого удара о брусок шарик останавливается, а брусок смещается по горизонтальной поверхности стола на расстояние $S$. Найдите:
	\lpu скорость бруска сразу после удара;\lpuend
	\lpu коэффициент трения скольжения между бруском и столом. \lpuend
	}
%
{\puan  $v = \sqrt{\frac{2gl}{7}};$~\puan $\mu = \frac{L}{49S}.$}
%%%%%%%%%%%%%%%%%%%%%%%%%%%%%%%%%%%%%%%%%%%%%%%%%%%% 
\z[{\begin{wrapfigure}[7]{r}{0.3\textwidth}%
		\vspace{-10pt}%
		\begin{flushright}%
 			\includegraphics[width=0.3\textwidth]{7.8.3}%
		\end{flushright}%
	\end{wrapfigure}%
	}]%
	{По горизонтальной поверхности пола движется со скоростью $u = 0.2$ м/с тележка со штативом, к которому на нити длиной $l = 0.25$ м привязан шар (см. рисунок). Пуля, летящая горизонтально со скоростью $59u$, попадает в шар и застревает в нём. Массы пули и шара $m$ и $11m$, масса тележки намного больше массы шара. Размеры шара малы по сравнению с длиной нити. Найдите:
	\lpu скорость шара $v_1$ относительно тележки сразу после попадания пули;\lpuend
	\lpu скорость шара $v_2$ относительно пола сразу после попадания пули;\lpuend
	\lpu максимальный угол отклонения шарика от вертикали. \lpuend
	}
%
{\puan  $v_1 = 5u = 1\,\text{м/с};$~\puan  $v_2 = 4u = 0.8\,\text{м/с};$~\puan $\alpha = \arccos\left(1 - \frac{25u^2}{2gl}\right) \approx 37\degree.$}
%%%%%%%%%%%%%%%%%%%%%%%%%%%%%%%%%%%%%%%%%%%%%%%%%%%%	
\z[{\begin{wrapfigure}[5]{r}{0.3\textwidth}%
		\vspace{-10pt}%
		\begin{flushright}%
 			\includegraphics[width=0.3\textwidth]{7.8.4}%
		\end{flushright}%
	\end{wrapfigure}%
	}]%
	{По гладкой горизонтальной поверхности стола движется со скоростью $u$ горка с неподвижной относительно горки шайбой на вершине горки (см. рисунок). Пуля, летящая горизонтально со скоростью $35u$, попадает в шайбу и застревает в ней. В результате шайба съезжает с горки, не отрываясь от её гладкой поверхности, и покидает горку. Массы пули и шайбы равны $m$ и $8m$, масса горки намного больше массы шайбы. Найдите:
	\lpu скорость шайбы $v_1$ относительно горки сразу после попадания пули; \lpuend
	\lpu скорость шайбы $v_2$ относительно стола сразу после попадания пули;\lpuend
	\lpu с какой скоростью относительно стола шайба покинула горку. \lpuend
	
	\noindent
	Известно, что при съезде с неподвижной горки изначально неподвижной шайбы шайба приобретает скорость $3u$.
		\newpage
	}
%
{\puan  $v_1 = 4u;$~\puan  $v_2 = 3u;$~\puan $v = 4u.$}
%%%%%%%%%%%%%%%%%%%%%%%%%%%%%%%%%%%%%%%%%%%%%%%%%%%% 
\z[{\begin{wrapfigure}[4]{r}{0.3\textwidth}%
		\vspace{-5pt}%
		\begin{flushright}%
 			\includegraphics[width=0.3\textwidth]{7.8.5}%
		\end{flushright}%
	\end{wrapfigure}%
	}]%
	{Доска 1 лежит на такой же доске 2. Обе они как целое скользят по гладкой ледяной поверхности со скоростью $v_0$ и сталкиваются с такой же доской 3, верхняя поверхность которой покрыта тонким слоем резины. При ударе доски 2 и 3 прочно сцепляются. Чему равна длина $l$ каждой доски, если известно, что доска 1 прекратила движение относительно досок 2 и 3 из-за трения после того, как она полностью переместилась с 2 на 3? Все доски твёрдые. Коэффициент трения между досками 1 и 3 равен $k$. Трением между досками 1 и 2, а также трением досок 2 и 3 о лёд можно пренебречь.
	}
%
{$l = \frac{v_0^2}{6kg}$}